% =============================================================================
%  CV de Emanuel Edgar Ancco Guaygua - Las Bambas 2026
%  Optimizado para: Seguridad de Presas (Vacante 503057)
% =============================================================================

\documentclass[10pt, a4paper]{article}

% --- PAQUETES NECESARIOS ---
\usepackage[utf8]{inputenc} 
\usepackage[spanish]{babel}      
\usepackage{geometry}            
\usepackage{hyperref}            
\usepackage{fontawesome5}        
\usepackage{titlesec}            
\usepackage{amsmath}             
\usepackage{lmodern} 

% --- CONFIGURACIÓN DE LA PÁGINA ---
\geometry{
    a4paper,
    left=1.6cm,
    right=1.6cm,
    top=1.3cm,
    bottom=1.3cm
}

% --- COLORES Y ENLACES ---
\hypersetup{
    colorlinks=true,
    linkcolor=black,
    urlcolor=blue,
    pdftitle={CV - Emanuel Edgar Ancco Guaygua - Las Bambas},
    pdfauthor={Emanuel Edgar Ancco Guaygua},
}

% --- PERSONALIZACIÓN DE TÍTULOS DE SECCIÓN ---
\titleformat{\section}
  {\Large\bfseries\scshape} 
  {}
  {0em}
  {}[\titlerule] 

\titlespacing*{\section}{0pt}{1.0ex plus 0.6ex minus .2ex}{0.6ex plus .2ex}

% --- ELIMINAR SANGRÍA DE PÁRRAFO ---
\setlength{\parindent}{0pt}

% --- DEFINICIÓN DE COMANDOS PARA ENTRADAS ---
\newcommand{\entry}[4]{
    \textbf{#1} \hfill \textit{#2} \\
    \textit{#3} \hfill \textit{#4} \\
    \vspace{-1.5mm} 
}

% =============================================================================
%  INICIO DEL DOCUMENTO
% =============================================================================
\begin{document}

% --- ENCABEZADO ---
\begin{center}
    {\Huge\bfseries EMANUEL EDGAR ANCCO GUAYGUA}
    \vspace{2mm}

    \faMapMarker*{} Lima / Puno, Perú \quad
    \faPhone*{} 974583549 \quad
    \faEnvelope{} \href{mailto:coarp.eancco@gmail.com}{coarp.eancco@gmail.com} \quad
    \faLinkedin{} \href{https://www.linkedin.com/in/emanuel-edgar-ancco-guaygua-a61210352}{LinkedIn}
    
    \faGlobe{} \textbf{Portafolio:} \href{https://emanuelancco.github.io/cv-portfolio/}{emanuelancco.github.io/cv-portfolio}
\end{center}

% --- RESUMEN PROFESIONAL ---
\section*{Resumen Profesional}
Egresado de Ingeniería Civil (UPC, Diciembre 2025) con experiencia en \textbf{gestión documentaria y elaboración de documentos técnicos} en obras públicas. Especialización en \textbf{análisis de datos con Python, R y MATLAB}, desarrollo de sistemas de Machine Learning para monitoreo estructural (arquitecturas PINN-GNN), y automatización de procesos con herramientas BIM. Conocimiento de estándares internacionales de seguridad de presas (\textbf{GISTM, ANCOLD, CDA}). Experiencia en \textbf{coordinación con equipos multidisciplinarios} y elaboración de informes técnicos. Originario de Puno, familiarizado con trabajo en altura.

\vspace{2mm}
\begin{center}
\fbox{\parbox{0.92\textwidth}{\centering
\faExternalLink*{} \textbf{EVIDENCIA TÉCNICA Y DEMOS:} \href{https://emanuelancco.github.io/cv-portfolio/}{emanuelancco.github.io/cv-portfolio} \\
\small{Videos de inferencia, métricas de modelos ML, arquitecturas GNN, capturas de plugins Revit}
}}
\end{center}

% --- EXPERIENCIA PREPROFESIONAL ---
\section*{Experiencia Preprofesional}

\entry{Asistente de Oficina Técnica}{Septiembre 2025 – Diciembre 2025}{Obra: ``Creación de Servicios Culturales - Teatro Municipal de Yunguyo''}{Puno, Perú}
\begin{itemize}
    \item Elaboración de informes mensuales de avance de obra, documentando progreso físico y financiero.
    \item Desarrollo de valorizaciones mensuales con cálculo de metrados ejecutados y control de costos.
    \item Modelado BIM estructural del proyecto (Revit 2026), generando modelos 3D para coordinación.
    \item Aplicación de plugins propios de automatización para armado de vigas y generación de losas.
\end{itemize}

\vspace{1mm} 

\entry{Asistente de Oficina Técnica}{Junio 2024 – Noviembre 2024}{Constructora Talgus S.A.C.}{Lima, Perú}
\begin{itemize}
    \item Elaboración de metrados y análisis de costos unitarios para preparación de presupuestos.
    \item Seguimiento y control de costos de obra, preparación de informes de valorización.
    \item Modelado y actualización de planos técnicos (AutoCAD, Revit) para correcta ejecución.
\end{itemize}

\vspace{1mm} 

\entry{Director del Área de Investigaciones}{Enero 2024 – Diciembre 2025}{Capítulo Estudiantil de Ingeniería Civil - IDI UPC}{Lima, Perú}
\begin{itemize}
    \item Gestión de proyectos de investigación aplicada en IA para infraestructura, liderando equipos multidisciplinarios.
    \item Supervisión del ciclo de vida de proyectos, desde planificación hasta presentación de resultados técnicos.
    \item Ponencia en CONEIC Ucayali 2024 sobre ``Detección Automatizada de Fisuras y Grietas con CNN''.
\end{itemize}

% --- EDUCACIÓN ---
\section*{Educación}
\entry{Pregrado en Ingeniería Civil (Egresado)}{2021 - Diciembre 2025}{Universidad Peruana de Ciencias Aplicadas (UPC)}{Lima, Perú}
\vspace{1mm}
\entry{Bachillerato Internacional (IB Diploma)}{2015 - 2017}{Colegio de Alto Rendimiento (COAR Puno)}{Puno, Perú}

% --- PROYECTOS DE DESARROLLO DE SOFTWARE ---
\section*{Proyectos de Desarrollo de Software e IA}

\textbf{GAIATECH SHM | Sistema de Monitoreo de Salud Estructural con IA}
\begin{itemize}
    \item Sistema de \textbf{detección de anomalías en puentes} usando datos de 5 acelerómetros y redes neuronales de grafos.
    \item Arquitectura innovadora \textbf{PINN-GNN} (Physics-Informed Neural Network) con grafo basado en geometría 3D real.
    \item Mejor modelo: \textbf{STG-AE Wavelet-GNN} con 5.1M parámetros y validation loss de \textbf{0.0064}.
    \item Extracción de features con \textbf{Discrete Wavelet Transform (DWT)} y análisis espectral (FFT, CWT).
    \item Tecnologías: Python, \textbf{PyTorch, PyTorch Geometric}, Wavelets (PyWavelets), CUDA.
\end{itemize}

\vspace{1mm}

\textbf{EMARC VISIÓN | Sistema de Detección de Grietas y EPP con Visión Artificial}
\begin{itemize}
    \item Sistema de \textbf{segmentación semántica de fisuras y grietas} en superficies de concreto (muros, pavimentos, puentes).
    \item Modelo \textbf{SegFormer} (Transformers) con clasificación automática de severidad: Leve/Moderada/Severa.
    \item Cálculo automático de métricas: \textbf{área (cm²), longitud, ancho promedio, porcentaje de daño}.
    \item Recomendaciones automáticas de mantenimiento según normativa (inspección/reparación).
    \item Módulo adicional: Detección de EPP con YOLOv8 (\textbf{mAP@0.5 = 90.7\%}).
    \item Tecnologías: Python, \textbf{Hugging Face Transformers}, SegFormer, YOLOv8, OpenCV.
\end{itemize}

\vspace{1mm}

\textbf{Plataforma de Predicción de Riesgos Laborales | NLP + Deep Learning}
\begin{itemize}
    \item API REST (FastAPI) para clasificación de narrativas de incidentes usando \textbf{DistilBERT}.
    \item Predicción de naturaleza del evento, parte del cuerpo afectada y tipo de evento.
    \item Módulo de \textbf{Toolbox Talks} automáticos para prevención proactiva de accidentes.
    \item Tecnologías: Python, \textbf{PyTorch, Transformers (Hugging Face)}, FastAPI, SQL.
\end{itemize}

\vspace{1mm}

\textbf{EMARC BIM SUITE | Suite de Automatización para Revit 2026}
\begin{itemize}
    \item \textbf{SAPtoRevit:} Importador de modelos SAP2000/ETABS a Revit (vigas, columnas, losas con pendientes).
    \item \textbf{EMARC Structural:} Suite de 20+ comandos para armado automático de elementos estructurales.
    \item \textbf{AutoCAD $\leftrightarrow$ Revit:} Interoperabilidad JSON para transferencia de geometría y familias.
    \item Tecnologías: C\#, \textbf{Revit API 2026}, .NET Framework 4.8.
\end{itemize}

\vspace{1mm}

\textbf{Innovación Tecnológica e Instrumentación IoT (Hardware Propio)}
\begin{itemize}
    \item \textbf{PachaGuard (IoT/FPGA):} Sistema de vigilancia autónoma con \textbf{FPGA Tang Nano 1K} y ESP32-CAM. Procesamiento en el borde (Edge Computing) para alertas en tiempo real vía UART.
    \item \textbf{Detección Ambiental (Aves/Tráfico):} Modelos \textbf{YOLOv11 y DETR} (Transformers) para monitoreo de fauna y flujo vehicular.
    \item \textbf{Análisis de Materiales:} Scripting avanzado con librería \textit{concrete-properties} (ACI) para diagramas momento-curvatura.
    \item Tecnologías: Verilog (FPGA), C++ (Arduino), MicroPython, Transformers.
\end{itemize}

\vspace{1mm}

\textbf{Herramientas de Análisis Estructural e Hidráulico}
\begin{itemize}
    \item \textbf{Análisis Sísmico Modal:} Herramienta Python para análisis dinámico con norma E.030 (matrices de rigidez, modos, derivas).
    \item \textbf{Simulación Monte Carlo:} Análisis de confiabilidad estructural con métodos FORM/FOSM y curvas de fragilidad.
    \item \textbf{Diseño Hidráulico:} Cálculo de canales con ecuación de Manning, exportación a Civil 3D, HEC-RAS, IBER.
    \item \textbf{Detección de Tráfico Vehicular:} Sistema YOLOv5/v11 para conteo y clasificación de vehículos en tiempo real.
    \item Tecnologías: Python, NumPy, SciPy, PyWavelets, OpenCV.
\end{itemize}

% --- HABILIDADES TÉCNICAS ---
\section*{Habilidades Técnicas}
\begin{itemize}
    \item \textbf{Análisis de Datos:} \textbf{Python, R, MATLAB} para procesamiento y visualización de datos, SQL, Power BI, Excel VBA.
    \item \textbf{Ofimática:} \textbf{MS Office nivel avanzado} (Word, Excel, PowerPoint), \textbf{SAP nivel básico}.
    \item \textbf{Programación y AI/ML:} Python (PyTorch, TensorFlow, Scikit-learn, YOLOv8), C\# (.NET).
    \item \textbf{Confiabilidad Estructural:} Métodos FORM/FOSM, Simulación Monte Carlo, Curvas de Fragilidad.
    \item \textbf{Ingeniería BIM/CAD:} Revit (Desarrollo de Plugins), AutoCAD, Civil 3D, SAP2000.
    \item \textbf{Automatización:} n8n (workflows), Dynamo, generación automática de reportes Word/Excel.
    \item \textbf{Sistemas GIS:} ArcGIS, QGIS.
    \item \textbf{Idiomas:} Español (Nativo), Inglés Intermedio-Avanzado (dominio técnico y conversacional).
\end{itemize}

% --- CERTIFICACIONES Y RECONOCIMIENTOS ---
\section*{Certificaciones y Reconocimientos}
\textit{\small Verificables en LinkedIn: \href{https://www.linkedin.com/in/emanuel-edgar-ancco-guaygua-a61210352}{linkedin.com/in/emanuel-edgar-ancco-guaygua}}
\vspace{1mm}
\begin{itemize}
    \item \textbf{Especialización en Machine Learning on Google Cloud} (4 cursos) | Google Cloud / Coursera, 2025
    \item \textbf{Especialista en IA y Aprendizaje Automático con Python} (120 horas) | Data Science Analysis, 2024
    \item \textbf{Ponente en CONEIC Ucayali 2024} - ``Aplicación de CNN para Evaluación Automatizada de Fisuras y Grietas en Construcciones de Lima''
    \item \textbf{2do lugar en Concurso de Conocimientos CONEIC} | Huaraz, 2025
    \item \textbf{Primer Puesto, Concurso de Ponencias Académicas} | XIV COREIC Lima 2023, UNFV
    \item \textbf{Innovador Destacado, DigiEduHack 2023} | Comisión Europea
    \item \textbf{Especialización en Python for Everybody} (5 cursos) | University of Michigan / Coursera
    \item \textbf{ISC AI Foundations For Everyone} | Campus Global ISC, 2025
    \item \textbf{ISC Data Analytics Aplicado a Cadena de Suministro} | Campus Global ISC, 2025
\end{itemize}

% --- INFORMACIÓN ADICIONAL ---
\section*{Información Adicional}
\begin{itemize}
    \item \textbf{Disponibilidad:} Inmediata para régimen atípico (14x7) en Unidad Minera.
    \item \textbf{Origen:} Puno, Perú - \textbf{Familiarizado con trabajo en altura} y condiciones de sierra (4,000+ msnm).
    \item \textbf{Conocimiento normativo:} Familiaridad con estándares GISTM, CDA y ANCOLD para seguridad de presas.
\end{itemize}

\end{document}
